\scp{Summary of tendencies observed in the region}{01-08}
In this section we discuss a number of tendencies that have been observed in this region,
some of which have been posited by others as subgrouping arguments,
and are discussed in the relevant chapters.
The diversity of the data in the \tsc{An} languages of eastern Indonesia
and Timor-Leste require reconstructing a proto language that is very close to PMP\footnote{
		Regarding the lexicon,
		\citet[245]{Blust1993} says,
		“Systematic attempts to reconstruct Swadesh 200-item test lists at
		various time-depths show clearly that Proto-\ili{Central}-Eastern Malayo-Polynesian
		was hardly distinct from Proto-Malayo-Polynesian (98 percent similar).”}
and one whose phonology
and lexicon are more like PMP than they are like EMP or \tsc{POc}.
 Many broad tendencies that have been noted may be valid as tendencies,
but are not exclusively shared innovations that are diagnostic for subgrouping,
since so many exceptions occur throughout the region
and the same patterns are also found widely in languages to the west
and north of the Blust line.
The list below is a sample drawn from
\citet{Blust1993,Blust2013,DonohueGrimes2008,GrimesCIP}, and \citet{Schapper2015}.

\begin{itemize}
	\item There is a tendency to lose PMP *y, *q, *h.
				But traces are documented in the region.
				Traces of *y and *q are particularly widespread
				and must be reconstructed,
				the latter in all positions.
				We revisit the issue of PMP *h in \srf{13-01-04}.
	\item There is a tendency towards a weakening of the
				pre-penultimate vowel in historical trisyllabic roots,
				particularly those that begin with *qa- and *ha-.
				But lexically specific traces are geographically widespread,
				and sometimes the historical antepenultimate syllable explains
				the difference between a reflex in word-initial position
				in contrast to a conditioned intervocalic reflex,
				so they must be reconstructed.
				This same tendency also occurs widely west of the Blust line.
	\item There is a tendency for historical vowel-glide sequences
				such as PMP \mbox{*-ay}, *-aw, *-uy, *-iw to reduce to a single vowel.
				This tendency also occurs widely west of the Blust line.
				Three strategies of reduction are common in the region,
				with a few languages also preserving the full vowel-glide sequence unreduced.
				For example: 
	\begin{itemize}
		\item the vowel (the first element) is retained
					and the glide is lost:
					*-ay > \it{a}, *-aw > \it{a},
					*uy > \it{u}, *-iw > \it{i}; 
		\item the glide (the second element)
					is retained as a vowel,
					and the historical vowel is lost:
					*{\nbhyp}ay > \it{i}, *-aw > \it{u}, *-uy > \it{i}, *-iw > \it{u}; 
		\item the vowel-glide sequence coalesces (merges,
					combines) to a single vowel that is different from
					the historical vowel: *-ay > \it{e}; *-aw > \it{o}.
	\end{itemize}
\end{itemize}

Consequently,
the whole historical vowel-glide sequence must be reconstructed for this region.
\citet{Blust1993} refers to the first pattern as “glide truncation”,
clarified in \citet[31]{Blust2009} as “truncation of the coda”
in all of the vowel-glide sequences.
“Glide truncation” is not found through the entire region,
and does not fully characterise any of our higher-level subgroups.
We revisit this issue in \srf{13-01-03}.

\begin{itemize}
		\item There is an areal tendency towards “postnasal voicing of stops”.
					\citet[72]{Blust2009} describes this as a “three-step change that involves:
					(1) loss of the vowel in PMP *ma- \tsc{‘stative’},
					(2) postnasal voicing of a base-initial voiceless stop,
					and (3) loss of the conditioning nasal.” 
	\begin{itemize}
			\item Examples of the postnasal voicing of stops
						would be PMP *ma-putiq > \ili{Buru} (\tsc{Sula-Buru}) \it{boti} ‘white’,
						*ma-panas ‘hot’ > \ili{Buru} \it{bana} ‘fire’.
			\item But we also note the presence of
						other patterns in the region,
						such as where the nasal is retained
						and the stop is lost, as in several \ili{Meto}
						(\tsc{Timor-Babar}) languages of west Timor,
						such as *ma-putiq > \it{mutiʔ} ‘white’
						and *ma-panas > \it{manas} ‘sun’.
			\item Or the vowel is lost from the historical prefix,
						such as in \ili{Kambera} (\tsc{\ili{Bima}-Lembata}) *ma-panas > \it{mbanah-u} ‘hot’.
			\item Or the nasal and labial merge but stay voiceless,
						as in \ili{Dhao} *ma-panas > \it{pana} ‘hot’
						(where commonly *p > {\0}, \it{p},
						but always *ma-p > \it{p}).
						This is also common in \tsc{\ili{Ambon}-Seram}
						languages of central Maluku.
			\item In some cases *ma-p is retained in full
						with no reduction and no voicing,
						such as \ili{Yamdena} \it{mafuti} ‘white’
						and \it{mefanas} ‘hot’, or \ili{Geser}-\ili{Gorom} \it{mafuti} ‘white’
						and \it{mahanas} ‘hot’
						(both \tsc{Seram-Tanimbar-Bomberai} languages).
	\end{itemize}
\end{itemize}

Blust’s “postnasal voicing of stops” does
not define any of our higher level subgroups,
but the sequencing of PMP *ma-p > **mp > \it{b}
becomes relevant in some lower order branches.
We revisit this issue in \srf{13-01-02}.

\begin{itemize}
\item There is a tendency to lose final consonants from lexical roots.
			But retentions are widespread,
			and most must be reconstructed (see \srf{07-02-01}).
			All subgroups have some retentions of some historical final consonants -- some more than others.
			Loss of some historical final consonants also occurs in whole subgroups west of the Blust line.
\item There is a geographically widespread lack of velar nasal /ŋ/,
			or a merger of PMP *ŋ and *n.
			But retentions are also widespread
			and contrast with nodes that merge *ŋ/*n > \it{n} which is relevant for subgrouping,
			so the distinction must be reconstructed.
			We revisit this issue in \srf{13-01-01}.
\item There is a tendency in the morphosyntax to have subject marking prefixes on some verbs;
			but this also occurs widely west of the Blust line.
			\citet[25]{Adelaar2005} has pointed out that the subject markers in “Moluccan languages”
			and “West Melanesian languages” are not reflected in cognate forms that are commonly inherited
			and could be from “convergent development,
			which potentially weakens the evidence for CEMP”.
			We note that the subject prefixes of \tsc{\ili{Ambon}-Seram} languages are also not cognate with the subject
			prefixes of the \tsc{Rote-Meto} languages (\tsc{Timor-Babar}).
			In fact, within the \tsc{Timor-Babar} subgroup itself,
			the subject prefixes of different branches are not all cognate with one another.
			The presence of subject markers does not define any of our higher level subgroups,
			but it is a feature of some lower level groups.
			This issue is revisited in several chapters.
\item The preposed possessor construction introduced in \srf{01-05}
			is an areal feature resulting from contact
			with typologically SOV Papuan languages,
			and should not be used for linguistic subgrouping.
			It is not found in Sumba,
			\ili{Havu}, \ili{Dhao}, nor in languages of western Flores.
			This issue is revisited in \srf{02-03-01-01}.
\item The sporadic presence of reflexes of {\#}muku ‘banana’\footnote{
				The hashtag on {\#}muku marks that this
				is a pseudo-etymon, and has reflexes in at l\ili{east} two unrelated language families.
				\citet{DonohueDenham2009} note evidence of several other etyma
				reconstructed for ‘banana’ for Austronesian languages in the region.
				Our discussion here focuses only on *punti
				and {\#}muku.}
			in languages from western Flores to
			\ili{Aru} and slightly further north to the Bomberai Peninsula,
			but not in the languages of central Maluku,
			is due to contact with Papuan languages,
			rather than genealogical in the linguistic sense.
			 It co-exists alongside reflexes of PMP *punti ‘banana’
			and several other etyma.
			This is discussed further in \srf{02-03-01-07},
			with fuller data in appendix \ref{ch:BanTer}.
\end{itemize}


To summarise, there are many observable tendencies in the Austronesian languages in the region,
yet none of those identified qualify as exclusively shared innovations,
nor do they meet criteria normally accepted for subgrouping arguments.
The fact that a feature is found sporadically throughout a region with much mobility
and trade (see \srf{02-02}),
and that this region has been called a “diffusion corridor” is not sufficient grounds for defining a subgroup.
We revisit several of these issues in chapters \ref{ch:OthSub}--\ref{ch:MorLes}.

We can confidently say that there are patterned differences that are loosely demarcated by the Blust
line, but none of the issues yet identified actually define the Blust line.
Many Austronesian languages \ili{east} of the Blust line display a number of features not found in languages
west of the line,
yet the ones described so far do not require or reflect a common parent language or linkage,
but very strongly reflect that this region has been a zone of intense contact with Papuan languages for
a very long time.
In fact, there are many lines,
or a cluster of different lines,
marking many different features (see \frf{15-12}
and distribution maps throughout this study).
There is not just one line
demarcating a particular feature or set of features.
So on the basis of what we have seen so far,
travelling from west to \ili{east} we should be thinking about entering a zone of transition
and contact, rather than crossing a narrow boundary in which there is a cohesive subgroup.
We are crossing a very wide threshold
and entering a mutually respectful
and accommodating meeting place,
not an Austronesian colony.

